\documentclass[12pt,openany]{book}

% ============= PACOTES
\usepackage[utf8]{inputenc}
\usepackage[T1]{fontenc}
\usepackage[brazil]{babel}

\usepackage{amsmath, amssymb, amsthm, mathtools}
\usepackage{enumitem}
\usepackage{geometry}
\usepackage{hyperref}
\usepackage{xcolor}
\usepackage[most]{tcolorbox}
\usepackage{titlesec}
\usepackage{fancyhdr}
\usepackage{graphicx}
\usepackage{tikz}
\usepackage{booktabs}
\usepackage{pgfplots}
\pgfplotsset{compat=1.18}

\geometry{margin=2.5cm}

% ============= COLORS
\definecolor{maincolor}{RGB}{0, 60, 120}
\definecolor{accentcolor}{RGB}{180, 30, 30}
\definecolor{lightcolor}{RGB}{230, 240, 255}
\definecolor{examplecolor}{RGB}{245, 245, 245}

% ============= LINKS
\hypersetup{
    colorlinks=true,
    linkcolor=maincolor,
    urlcolor=maincolor
}

% ============= PARAGRAPHS
\setlength{\parindent}{0pt}
\setlength{\parskip}{0.6em}

% ============= HELPERS
\newcommand{\R}{\mathbb{R}}
\newcommand{\N}{\mathbb{N}}
\newcommand{\Z}{\mathbb{Z}}
\newcommand{\Q}{\mathbb{Q}}
\newcommand{\eps}{\varepsilon}

\DeclareMathOperator{\im}{im}
\DeclareMathOperator{\dom}{dom}

% ============= CHAPTER STYLE
\titleformat{\chapter}[display]
  {\Huge\bfseries\color{maincolor}}
  {\filleft\Huge\thechapter}
  {2ex}
  {\titlerule\vspace{1ex}\filright}
  [\vspace{1ex}\titlerule]

% ============= HEADER/FOOTER
\pagestyle{fancy}
\fancyhf{}

\fancyhead[L]{\nouppercase{\leftmark}}
\fancyhead[R]{$\sum$ymbolica}

\fancyfoot[C]{\thepage}

\renewcommand{\headrulewidth}{0.4pt}
\renewcommand{\footrulewidth}{0pt}

% ============= BOX (tcolorbox)
\tcbset{
    defstyle/.style={
        colback=lightcolor,
        colframe=maincolor,
        fonttitle=\bfseries,
        sharp corners
    },
    thmstyle/.style={
        colback=green!5,
        colframe=green!50!black,
        fonttitle=\bfseries,
        sharp corners
    },
    examplestyle/.style={
        colback=examplecolor,
        colframe=black!50,
        fonttitle=\bfseries,
        sharp corners
    },
    intuitionstyle/.style={
        colback=yellow!10,
        colframe=yellow!50!black,
        fonttitle=\bfseries,
        sharp corners
    },
    exercise/.style={
        colback=red!5,
        colframe=red!60!black,
        fonttitle=\bfseries,
        sharp corners
    }
}

% ============= ENVIRONMENTS
\newtcbtheorem[number within=chapter]{definition}{Definição}{defstyle}{def}
\newtcbtheorem[number within=chapter]{theorem}{Teorema}{thmstyle}{th}
\newtcbtheorem[number within=chapter]{lemma}{Lema}{thmstyle}{lem}
\newtcbtheorem[number within=chapter]{corollary}{Corolário}{thmstyle}{cor}
\newtcbtheorem[number within=chapter]{proposition}{Proposição}{thmstyle}{prop}
\newtcbtheorem[number within=chapter]{example}{Exemplo}{examplestyle}{ex}
\newtcbtheorem[number within=chapter]{remark}{Observação}{examplestyle}{rem}
\newtcbtheorem[number within=chapter]{exercise}{Exercício}{exercise}{exe}

\newtcolorbox{intuitionbox}{
    intuitionstyle,
    title=Intuição
}

\renewcommand{\proofname}{Prova}

% ============= COVER
\newcommand{\makebookcover}{
\begin{titlepage}
    \begin{tikzpicture}[remember picture, overlay]
        \fill[maincolor] (current page.north west)
            rectangle ([xshift=4cm]current page.south west);
    \end{tikzpicture}

    \vspace*{3cm}

    \hspace{4.5cm}
    \begin{minipage}{0.6\textwidth}
        {\Huge\bfseries \textcolor{maincolor}{Métodos Numéricos para Equações Diferenciais} \par}
        \vspace{0.5cm}
        {\Large Notas de Estudo \par}

        \vspace{2cm}

        {\Large $\sum$ymbolica \par}

        \vfill

        {\large \today \par}
    \end{minipage}
\end{titlepage}
}

\usepackage{tikz}

% ============= DOCUMENT
\begin{document}

\makebookcover

\tableofcontents

\chapter{Fundamentos}

\section{Derivada Numérica}
\subsection{Diferenças Finitas}
Uma vez que a definição de $f'(x)$ é dada por
\[
    f'(x) = \lim_{\Delta x \to 0} \frac{f(x_0) + \Delta x - f(x_0)}{\Delta x}
\]

podemos aproximá-la \textbf{numericamente} num ponto qualquer de seu domínio a partir da reta secante entre dois pontos.

\begin{enumerate}[label=(\roman*)]
    \item \textbf{Diferença Finita Progressiva}:
        \[
            f'(x_0) \approx \frac{f(x_0 + h) - f(x_0)}{h}
        \]
    \item \textbf{Diferença Finita Regressiva}:
        \[
            f'(x_0) \approx \frac{f(x_0) - f(x_0 + h)}{h}
        \]
    \item \textbf{Diferença Finita Centrada}:
        \[
            f'(x_0) \approx \frac{f(x_0 + h) - f(x_0 - h)}{2h}
        \]
\end{enumerate}

\begin{remark}{Comentário}{}
    Note que podemos escrever a diferença finita centrada como a média da diferença finita progressiva e regressiva.
\end{remark}

\begin{example}{}{}

Aproxime a derivada de $f(x) = \ln x$ em $x_0 = 2$.

\textbf{Caso 1: $h = 0.1$}

\begin{enumerate}[label=(\roman*)]
    \item Diferença progressiva:
    \[
        f'(2) \approx \frac{\ln 2.1 - \ln 2}{0.1} \approx 0.487902
    \]

    \item Diferença regressiva:
    \[
        f'(2) \approx \frac{\ln 2 - \ln 1.9}{0.1} \approx 0.512933
    \]

    \item Diferença centrada:
    \[
        f'(2) \approx \frac{\ln 2.1 - \ln 1.9}{2 \cdot 0.1} \approx 0.500417
    \]
\end{enumerate}

\vspace{1em}

\textbf{Caso 2: $h = 0.01$}

\begin{enumerate}[label=(\roman*)]
    \item Diferença progressiva:
    \[
        f'(2) \approx \frac{\ln 2.01 - \ln 2}{0.01} \approx 0.498754
    \]

    \item Diferença regressiva:
    \[
        f'(2) \approx \frac{\ln 2 - \ln 1.99}{0.01} \approx 0.501254
    \]

    \item Diferença centrada:
    \[
        f'(2) \approx \frac{\ln 2.01 - \ln 1.99}{2 \cdot 0.01} \approx 0.500004
    \]
\end{enumerate}

\end{example}

Uma outra forma de visualizar as diferenças divididas é através de \textbf{Operadores de Diferenças}.
A diferença dividida do \textbf{Método de Newton} de dois pontos é uma diferença finita:
\[
    D_+f(x) = \frac{f(x + h) - f(x)}{h}
\]
\[
    D_-f(x) = \frac{f(x) - f(x-h)}{h},
\]
\[
    D_of(x) = \frac{f(x + h) - f(x - h)}{2h}
\]

\begin{remark}{Comentário}{}
    Note que
    \[
        D_of(x) = \frac{1}{2}\left[ D_+f(x) - D_-f(x) \right]
    \]
\end{remark}

Podemos aplicar sucessivos operadores para obter derivadas de ordem mais alta.
Aplicando $D_+$ e em seguida $D_-$ à função $f$ teremos
\[
    \begin{aligned}
        D^2f(x) &= D_-\left( D_+f(x) \right) \\
        &= D_-\left( \frac{f(x + h) - f(x)}{h} \right) \\
        &= \frac{1}{h}\left( D_-f(x + h) - D_-f(x) \right) \\
        &= \frac{1}{j}\left[ \left(\frac{f(x + h) - f(x)}{h}\right) - \left( \frac{f(x) - f(x + h)}{h} \right) \right] \\
        &= \frac{1}{h^2}\left[ f(x + h) - 2f(x) + f(x - h) \right]
    \end{aligned}
\]

Temos então que
\[
    f''(x) \approx \frac{1}{h^2}\left[ f(x + h) - 2f(x) + f(x - h) \right]
\]

\begin{remark}{Comentário}{}
    Essa é a chamada diferença finita centrada para a 2º derivada.
\end{remark}

Valem também as seguintes igualdades
\[
    D_+[D_-f(x)] = D_-[D_+f(X)] = D^2_o\left[ \frac{h}{2} \right]f(x)
\]

em que $D^2_o\left[\frac{h}{2}\right]$ representa a diferença dividida centrada a um passo de $\frac{h}{2}$.
A prova dessa igualdade está disponível no material complementar.

Por fim, podemos ainda obter as diferenças divididas através da \textbf{Série De Taylor}.
Abaixo, a série para a função $f$ ao redor do ponto $x_0$:
\[
f(x) = \sum_{n=0}^{\infty} f^{(n)}(x_0)\frac{(x - x_0)^n}{n!}
\]

Para um passo à frente, tomando $x - x_0 = h$, temos
\begin{equation}\label{cap_1:progressive_division_difference}
    f(x_0 + h) \approx f(x_0) + f'(x_0)h + f''(x_0)\frac{h^2}{2} + f'''(x_0)\frac{h^3}{3!}.
\end{equation}

Observe que, parando na primeira ordem, obtemos a diferença dividida progressiva
\[
    f(x_0 + h) \approx f(x_0) + f'(x_0)h \Rightarrow f'(x_0) \approx \frac{f(x_0 + h) - f(x_0)}{h}.
\]

Para um passo atrás, tomando $x - x_0 = -h$, temos
\begin{equation}\label{cap_1:regressive_division_difference}
    f(x_0 - h) \approx f(x_0) - f'(x_0)h + f''(x_0)\frac{h^2}{2} - f'''(x_0)\frac{h^3}{3!}.
\end{equation}

Observe que, parando na primeira ordem, obtemos a diferença dividida regressiva
\[
    f'(x_0) \approx \frac{f(x_0) - f(x_0 - h)}{h}
\]

Podemos também subtrair \ref{cap_1:regressive_division_difference} de \ref{cap_1:progressive_division_difference} até segunda ordem, obtendo
\[
    f'(x_0) \approx \frac{f(x_0 + h) - f(x_0 - h)}{2h}.
\]

Somando-as até segunda ordem, obtemos
\[
    f''(x_0) \approx \frac{f(x_0 + h) - 2f(x_0) + f(x_0 - h)}{h^2}.
\]

As demais derivadas podem ser obtidas de forma similares.
Abaixo, são apresentadas as fórmulas gerais para as diferenças finitas:

\begin{enumerate}[label=(\roman*)]
    \item \textbf{Diferença Progressiva}:
        \[
            f^{(n)}(x) \approx \frac{1}{h^{n}}\left[ \sum_{k = 0}^{n} (-1)^{n - k} \binom{n}{k} f(x + kh) \right]
        \]
    \item \textbf{Diferença Regressiva}:
        \[
            f^{(n)}(x) \approx \frac{1}{h^{n}}\left[ \sum_{k = 0}^{n} (-1)^{k} \binom{n}{k} f(x - kh) \right]
        \]
    \item \textbf{Diferença Centrada}:
        \[
            f^{(n)}(x) \approx \frac{1}{h^{n}}\left[ \sum_{k = 0}^{n} (-1)^{k} f\left( x + \left(\frac{n}{2} - k \right)h \right) \right]
        \]
\end{enumerate}

\begin{example}{}{}
    Aproxime $f^{(3)}(x)$ para $f(x) = \sin(x)$ em $x_0 = 2$ para $h = 0.01$.

    \vspace{1em}

    \textbf{Solução:}
    \begin{enumerate}[label=(\roman*)]
        \item \textbf{Diferença progressiva}:
        
            Note que, pela fórmula geral, temos:
            \[
                f^{(3)}(x) \approx \frac{1}{h^{3}}\left[ -f(x) + 3f(x + h) - 3f(x + 2h) + f(x + 3h) \right].
            \]
            Portanto:
            \[
                \begin{aligned}
                    \sin^{(3)}(x) &\approx \frac{1}{0.01^3}\left[ -\sin(2) + 3\sin(2.01) - 3\sin(2.02) + \sin(2.03) \right] \\
                    &\approx 0.42973359937
                \end{aligned}
            \]
        
        \item \textbf{Diferença regressiva}:
            
            Note que, pela fórmula geral, temos:
            \[
                f^{(3)}(x) \approx \frac{1}{h^{3}}\left[ f(x) - 3f(x + h) + 3f(x + 2h) - f(x + 3h) \right].
            \]
            Portanto:
            \[
                \begin{aligned}
                    \sin^{(3)}(x) &\approx \frac{1}{0.01^3}\left[ \sin(2) - 3\sin(2.01) + 3\sin(2.02) - \sin(2.03) \right] \\
                    &\approx 0.40245604027
                \end{aligned}
            \]

        \item \textbf{Diferença centrada}:
        
            Note que, pela fórmula geral, temos:
            \[
                f^{(3)}(x) \approx \frac{1}{h^{3}}\left[ f\left(x + \frac{3}{2}h\right) - 3f\left(x + \frac{1}{2}h\right) + 3f\left(x + \frac{1}{2}h\right) \right].
            \]
            Portanto:
            \[
                \begin{aligned}
                    \sin^{(3)}(x) &\approx \frac{1}{0.01^3}\left[ \sin(2.015) - 3\sin(2.005) + 3\sin(1.995) - \sin(1.985) \right] \\
                    &\approx 0.41614163437
                \end{aligned}
            \]
    \end{enumerate}

    Lembrando que
    \[
        \sin^{(3)}(2) = -\cos(2) \approx 0.41614683655
    \]
\end{example}

Para derivadas de ordem ímpar podemos ainda construir fórmulas de diferença finitas centradas com passos com coeficientes inteiros.
Tomemos a média entre $f^{(n)}\left(x + \frac{h}{2}\right)$ e $f^{(n)}\left(x - \frac{h}{2}{}\right)$:
\[
    ...
\]

\begin{example}{}{}
    Para a derivada terceira da $f$ a um passo h, temos
    \[
        \begin{aligned}
        f^{(3)}(x)
        &\approx \frac{1}{2}\left[
            f^{(3)}\left(x + \frac{h}{2}\right)
            + f^{(3)}\left(x - \frac{h}{2}\right)
        \right]
        \\[6pt]
        &\approx \frac{1}{2h^{3}}
        \Bigl[
            \bigl(
                f(x + 2h)
                - 3f(x + h)
                + 3f(x)
                - f(x - h)
            \bigr)
        \\
        &\qquad\qquad
            +
            \bigl(
                f(x + h)
                - 3f(x)
                + 3f(x - h)
                - f(x - 2h)
            \bigr)
        \Bigr]
        \\[6pt]
        &= \frac{1}{2h^{3}}
        \left[
            f(x + 2h)
            - 2f(x + h)
            + 2f(x - h)
            - f(x - 2h)
        \right]
        \end{aligned}
    \]

    Aproximando, novamente, $\sin^{(3)}(2)$, agora através da fórmula apresentada, obtemos
    \[
        \begin{aligned}
            \sin^{(3)}(2) &\approx \frac{1}{2 \cdot 0.01^3} \left[ \sin(2.02) - 2\sin(2.01) + 2\sin(1.99) - \sin(1.98) \right] \\
            &\approx 0.41613643292
        \end{aligned}
    \]
\end{example}

\subsection{Diferenças Finitas a partir de Dados Tabelados}
Seja a seguinte tabela de dados:

\begin{table}[h]
    \centering
    \begin{tabular}{cc}
        \toprule
        $x_i$ & $y_i$ \\
        \midrule
        $x_1$ & $y_1$ \\
        $x_2$ & $y_2$ \\
        $\vdots$ & $\vdots$ \\
        $x_n$ & $y_n$ \\
        \bottomrule
    \end{tabular}
\end{table}

podemos aplicar, quando $x_{i + 1} - x_i = h, \forall i = 1, 2, \dots, n - 1$, os conhecimentos obtidos no capítulo.
Observada a restrição de que, para $y'_1$, deve-se utilizar a aproximação progressiva e, para $y'_n$, a regressiva.

\begin{example}{}{}
    Dada a tabela abaixo, aproxime $y'_i$ através dos métodos aprendidos.

    \begin{center}
        \begin{tabular}{cc}
            \toprule
            $x_i$ & $y_i$ \\
            \midrule
            $2.1$ & $0.7419$ \\
            $2.2$ & $0.7885$ \\
            $2.3$ & $0.8325$ \\
            $2.4$ & $0.8755$ \\
            $2.5$ & $0.9163$ \\
            \bottomrule
        \end{tabular}
    \end{center}

    \vspace{1em}

    \textbf{Solução:}
    A construção da tabela abaixo segue as fórmulas anteriores.
    Cabe o leitor verificar a corretude.

    \begin{center}
        \begin{tabular}{ccc}
            \toprule
            $x_i$ & $y_i$ & $y'_i$ \\
            \midrule
            $2.1$ & $0.7419$ & $0.466$ \\
            $2.2$ & $0.7885$ & $0.458$ \\
            $2.3$ & $0.8325$ & $0.435$ \\
            $2.4$ & $0.8755$ & $0.417$ \\
            $2.5$ & $0.9163$ & $0.408$ \\
            \bottomrule
        \end{tabular}
    \end{center}

    \textbf{Observação:} Note que, agora, uma vez obtidos $y'_i$, podemos utilizá-los para obtermos $y''_i$ e, de forma iterativa, derivadas de ordem superiores.
\end{example}

\subsection{Melhorando a Precisão}
Uma vez que a primeira derivada faz uso de apenas dois pontos, é necessário tomar um passo cada vez menor.
No entanto, o esforço computacional pode, por vezes, explodir rapidamente, como será visto mais adiante.
Dito isso, podemos aproximar a primeira derivada através da Série de Taylor truncada em segunda ordem:
\[
    f(x + h) \approx f(x) + f'(x)h + f''(x)\frac{h^{2}}{2}
\]
no entanto, sabemos que
\[
    f''(x) \approx \frac{1}{h^{2}}\Bigl[ f(x + h) - 2f(x) + f(x - h) \Bigr].
\]
Portanto,
\[
    f'(x) \approx \frac{1}{2h}\Bigl[ -3f(x) + 4f(x + h) - f(x + 2h) \Bigr]
\]
para a progressiva, e, para a regressiva,
\[
    f'(x) \approx \frac{1}{2h}\Bigl[ 3f(x) - 4f(x - h) + f(x - 2) \Bigr]
\]

Isso é particularmente útil para dados tabelados, uma vez que a diferença finita regressiva e progressiva apresentada anteriormente é de primeira ordem.

\end{document}